Since the SDP~2 hardware has not yet arrived, we built and tested everything that does not depend on it. This chapter reports what exists today, how it was checked and what was measured. Every number here comes from a run on a development PC rather than on the Raspberry Pi 5, and we keep the two apart throughout, leaving the Pi results for SDP~2.

\section{What has been built}
\label{sec:built}

The full software of Chapter 4 now exists as one repository of 86 source files, with about 5500 lines of Python, 1300 lines of web code and one small C++ library, running from a camera or a video file through to the web dashboard. Table \ref{tab:built} lists the parts. The modules follow the class diagram of Figure \ref{fig:ch4-10}, and the detector backend is selected in one place, so the same pipeline runs on PyTorch, \gls{onnx}, NCNN or, once the AI HAT+ arrives, the Hailo runtime.

\begin{table}[htbp]
\centering
\caption{Software built so far}
\label{tab:built}
\footnotesize
\setlength{\tabcolsep}{3pt}
\renewcommand{\arraystretch}{1.15}
\begin{tabular}{|p{0.190\linewidth}|p{0.500\linewidth}|p{0.240\linewidth}|}
\hline
\textbf{Part} & \textbf{What it contains} & \textbf{Technology} \\
\hline
Edge pipeline & Capture, crop, blind check, detector, \gls{ema}, decision state machine, alerts, calibration, eye module, crop tracking & Python, C++, OpenCV, Ultralytics, MediaPipe, YuNet \\
\hline
Event log & Driver, session, event, calibration and tip tables, integrity check at boot, recovery of unfinished sessions & SQLite in \gls{wal} mode \\
\hline
Driver display & \gls{pin} login, calibration, pre-trip check, status, settings, rest-stop card & Flask, HTML, 1024$\times$600 \\
\hline
Alert assets & 3 alarm sounds from \gls{matlab}, 40 voice messages in English and Arabic (80 clips) & WAV files, neural text-to-speech \\
\hline
Trip assistant & Nearest rest stop from OpenStreetMap, with a saved list of 363 places as a fallback & Overpass \gls{api} \\
\hline
Upload and server & \gls{tls} 1.3 uploader with device key, retry and duplicate handling, login, statistics, tips & FastAPI, PostgreSQL or SQLite \\
\hline
Dashboard & Overview, trips, statistics, tips and privacy pages, with hand-drawn charts & React, Vite \\
\hline
Dataset tools & Recording guide, frame extraction, subject-disjoint folds, annotator agreement, hard-negative mining & Python \\
\hline
Model tools & Training, evaluation, structured pruning, \gls{onnx} and NCNN export, Hailo compile script, Pi benchmark & Ultralytics, torch-pruning \\
\hline
Deployment & Pi setup script, boot service, kiosk start & Bash, systemd \\
\hline
\end{tabular}
\end{table}

The detector in every run below is our \gls{yolo}11n trained on DMD in the benchmark of Chapter 4, rather than a model trained on our own subjects, since that dataset has not yet been recorded.

\section{How it was tested}

Testing relied on two kinds of evidence: automated scripts that check the logic against the test plan of Section 5.4, and a replay that exercises the whole pipeline on real footage. The replay covers 648\,s of held-out DMD driver \texttt{subject\_05}. Since the DMD frames we hold are sampled at 1 \gls{fps}, each frame was repeated ten times to form a 10 \gls{fps} clip with a matching label file. This suffices to exercise the full system, although event times on this clip are only accurate to about a second.

All 194 automated checks pass across the eight test scripts of Table \ref{tab:checks}. These are software checks, where the light bar and audio are replaced by recording stand-ins and power cuts are simulated by killing the writing process.

\begin{table}[htbp]
\centering
\caption{Automated software checks}
\label{tab:checks}
\footnotesize
\setlength{\tabcolsep}{3pt}
\renewcommand{\arraystretch}{1.15}
\begin{tabular}{|p{0.210\linewidth}|p{0.610\linewidth}|p{0.100\linewidth}|}
\hline
\textbf{Script} & \textbf{What it covers} & \textbf{Passed} \\
\hline
Decision & \gls{ema} at 10, 15 and 30 \gls{fps}, confirmation, cool-down, tiers, eye rules & 29/29 \\
\hline
Logger & Schema, \gls{pin} storage, 50 process kills, integrity recovery & 16/16 \\
\hline
Vision & Detector, crop, blind check on synthetic and real frames, C++ against OpenCV & 25/25 \\
\hline
Alerts & Flash limit, sound and voice selection, clip coverage, muted cues & 14/14 \\
\hline
Eye and calibration & Calibration formula, \gls{ear}, \gls{perclos}, closure timing, look-down rules, crop tracking on a real face & 42/42 \\
\hline
Display & \gls{pin} flow, calibration, pre-trip check, display states & 26/26 \\
\hline
Upload & \gls{tls} refusal of bad certificates, device key, retry, duplicates, per-driver access, tips & 29/29 \\
\hline
Tools & Subject-disjoint folds, agreement measure, dataset helpers & 13/13 \\
\hline
\textbf{Total} & & \textbf{194/194} \\
\hline
\end{tabular}
\end{table}

\section{Results so far}
\subsection{Full pipeline replay}

Figure \ref{fig:replay} and Table \ref{tab:replay} present the replay. All four phone cues were caught and escalated, with tier times landing exactly at 5\,s and 10\,s after the cutoff crossing. Two of the four drinking cues were missed, at 148\,s and 227\,s, both detector misses on short cues from a driver the model never saw, and no detector alert fell outside a labelled cue. One result points to a design issue rather than a model issue: during the last long phone call, a brief drop in detection counted as the end of the cue, so the rest of the call fell inside the 30\,s cool-down and stayed silent.

\begin{figure}[htbp]
\centering
\begin{tikzpicture}[font=\footnotesize, x=1cm, y=0.72cm]
\draw[gray!25] (0.000,0.1) -- (0.000,4.5);
\node[anchor=north, font=\scriptsize] at (0.000,0.1) {0};
\draw[gray!25] (2.006,0.1) -- (2.006,4.5);
\node[anchor=north, font=\scriptsize] at (2.006,0.1) {100};
\draw[gray!25] (4.012,0.1) -- (4.012,4.5);
\node[anchor=north, font=\scriptsize] at (4.012,0.1) {200};
\draw[gray!25] (6.019,0.1) -- (6.019,4.5);
\node[anchor=north, font=\scriptsize] at (6.019,0.1) {300};
\draw[gray!25] (8.025,0.1) -- (8.025,4.5);
\node[anchor=north, font=\scriptsize] at (8.025,0.1) {400};
\draw[gray!25] (10.031,0.1) -- (10.031,4.5);
\node[anchor=north, font=\scriptsize] at (10.031,0.1) {500};
\draw[gray!25] (12.037,0.1) -- (12.037,4.5);
\node[anchor=north, font=\scriptsize] at (12.037,0.1) {600};
\node[font=\scriptsize] at (6.5,-0.85) {Replay time (s)};
\node[anchor=east] at (-0.1,4.0) {Labelled cues};
\node[anchor=east] at (-0.1,3.0) {Tier 1};
\node[anchor=east] at (-0.1,2.2) {Tier 2};
\node[anchor=east] at (-0.1,1.4) {Tier 3};
\node[anchor=east] at (-0.1,0.6) {Muted};
\fill[out1] (0.943,3.82) rectangle (1.003,4.18);
\fill[out1] (2.026,3.82) rectangle (2.147,4.18);
\fill[out1] (2.969,3.82) rectangle (3.049,4.18);
\fill[out1] (4.554,3.82) rectangle (4.634,4.18);
\fill[own1] (5.637,3.82) rectangle (6.219,4.18);
\fill[own1] (7.162,3.82) rectangle (7.844,4.18);
\fill[own1] (9.228,3.82) rectangle (10.151,4.18);
\fill[own1] (11.154,3.82) rectangle (11.937,4.18);
\fill[out1] (0.983,3.0) circle (1.1mm);
\fill[out1] (2.127,3.0) circle (1.1mm);
\fill[own1] (5.698,3.0) circle (1.1mm);
\fill[own1] (5.758,2.2) circle (1.1mm);
\fill[own1] (5.858,1.4) circle (1.1mm);
\fill[own1] (7.222,3.0) circle (1.1mm);
\fill[own1] (7.282,2.2) circle (1.1mm);
\fill[own1] (7.383,1.4) circle (1.1mm);
\fill[own1] (9.289,2.2) circle (1.1mm);
\fill[own1] (9.449,1.4) circle (1.1mm);
\fill[oth1] (10.853,1.60) -- (10.743,1.26) -- (10.963,1.26) -- cycle;
\fill[own1] (11.235,2.2) circle (1.1mm);
\draw[own1, line width=0.8pt, fill=white] (11.495,0.6) circle (1.1mm);
\draw[own1, line width=0.8pt, fill=white] (11.696,0.6) circle (1.1mm);
\draw[own1, line width=0.8pt, fill=white] (11.856,0.6) circle (1.1mm);
\fill[out1] (1.2,-1.62) circle (1.1mm); \node[anchor=west] at (1.32,-1.62) {Drinking};
\fill[own1] (3.6,-1.62) circle (1.1mm); \node[anchor=west] at (3.72,-1.62) {Phone use};
\fill[oth1] (6.2,-1.48) -- (6.09,-1.76) -- (6.31,-1.76) -- cycle; \node[anchor=west] at (6.42,-1.62) {Eye alarm};
\draw[gray, line width=0.8pt] (8.8,-1.62) circle (1.1mm); \node[anchor=west] at (8.92,-1.62) {Muted, logged only};
\end{tikzpicture}

\caption{Replay of held-out DMD driver \texttt{subject\_05}: labelled cues against the alerts the system raised}
\label{fig:replay}
\end{figure}

\begin{table}[htbp]
\centering
\caption{Replay results on PC hardware}
\label{tab:replay}
\footnotesize
\setlength{\tabcolsep}{3pt}
\renewcommand{\arraystretch}{1.15}
\begin{tabular}{|p{0.430\linewidth}|p{0.220\linewidth}|p{0.250\linewidth}|}
\hline
\textbf{Measure} & \textbf{Result} & \textbf{Target} \\
\hline
Phone cues caught & 4 of 4 & Recall $\ge$ 0.90 \\
\hline
Drinking cues caught & 2 of 4 & Recall $\ge$ 0.90 \\
\hline
Tier 2 and Tier 3 timing & 5.0\,s and 10.0\,s & $\pm$0.2\,s \\
\hline
Detector alerts outside labelled cues & 0 in 648\,s & $\le$ 3 per hour \\
\hline
Eye alarms on the same clip & 1 (26 with the first rule) & $\le$ 3 per hour \\
\hline
Latency, RTX 4060, no eye module & 37 ms & Pi target: $\le$ 100 ms \\
\hline
Latency, CPU, no eye module & 111.4 ms mean, 116.1 ms p95 & Pi target: $\le$ 100 ms \\
\hline
Latency, CPU, with eye module & 143.1 ms mean, 181.0 ms p95 & Pi target: $\le$ 100 ms \\
\hline
\end{tabular}
\end{table}

These numbers serve as a working check rather than an evaluation. One driver and eight cues cannot estimate recall, and 648\,s cannot estimate a false alarm rate, so the final figures will come from the leave-one-subject-out test of Section 5.4. The PC latencies also say nothing about the Pi, where detection runs on the \gls{npu}, and they overstate the work done per frame, since a replay keeps every frame and the figures include the time a frame waits in a queue. Timed directly, the detector and eye module together take about 71 ms per frame on the CPU.

\subsection{Logging, upload, compression and other checks}

Table \ref{tab:other-results} lists the other results.

\begin{table}[htbp]
\centering
\caption{Other results on the PC}
\label{tab:other-results}
\footnotesize
\setlength{\tabcolsep}{3pt}
\renewcommand{\arraystretch}{1.15}
\begin{tabular}{|p{0.150\linewidth}|p{0.480\linewidth}|p{0.300\linewidth}|}
\hline
\textbf{Item} & \textbf{Result} & \textbf{Still open} \\
\hline
Logging & 50 process kills left the database intact, with no committed event lost & Real relay test, because \texttt{synchronous=NORMAL} can lose the last commits in a power cut \\
\hline
Upload and dashboard & 41 replay events uploaded over \gls{tls} 1.3 and shown in the dashboard. Bad certificates refused, no duplicates, each login sees only its own data & PostgreSQL, Docker and hosting not run \\
\hline
Pruning & A 20\% target removed 24.9\% of the parameters and dropped \gls{map}$_{50}$ to 0 until fine-tuning, and 8 epochs recovered it & Full comparison with the acceptance rule \\
\hline
Export & \gls{onnx} and NCNN match PyTorch (confidence 0.818) & FP16, \gls{int}, \gls{qat} and the Hailo compile \\
\hline
Trip assistant & The saved list of 363 places answered in 4.31\,s after the live query timed out at 4\,s & A measured live answer \\
\hline
Blind check & Brightness 59 to 107 and sharpness 126 to 365 on DMD frames, with no false ``camera blocked'' & Limits tuned on the C922 \\
\hline
\end{tabular}
\end{table}

\subsection{Eye closure}
\label{sec:eye-results}

The first eye rule, an \gls{ear} below 75\% of the open-eye value for 1.5\,s, did not survive real footage. Across all 15{,}723 frames of the 14 alert DMD drivers, it fired about 100 times per hour against a target of 3. In a sample of 192 flagged frames, about a third showed a visible iris, during a smile, a squint or a slight downward look, while about 40\% were a real lid closure with the head looking down at the lap or the phone. From a dashboard camera, a downward look and a closed eye can appear identical in a single frame.


The rule now has three changes, all reflected in Table \ref{tab:ch4-15}. The \gls{ear} must fall below 50\% of the open-eye value, MediaPipe's eye-blink score must also exceed 0.50, and when the head is more than 8$^\circ$ below the driver's rest pitch, the alarm waits 3\,s instead of 1.5\,s. This last step only delays the alarm and never mutes it, as a driver nodding off also has the head down. Head-down time is also excluded from \gls{perclos}, and the rest pitch is measured during calibration alongside the open-eye \gls{ear}.

We tuned these values on drivers 1 to 9 and checked them on drivers 10 to 14. Figure \ref{fig:eye-rate} shows a drop of about 80\% that holds on the unseen drivers. At 1 \gls{fps} the rate can only be bounded, so the figure counts runs of two frames, which overcount, and runs of three frames, which undercount, and even the lower figure remains above the target of 3 per hour. On the replay clip, the eye alarms fell from 26 to 1.

\begin{figure}[htbp]
\centering
\begin{tikzpicture}
\begin{axis}[width=0.92\linewidth, height=6.2cm, ybar, bar width=20pt, area legend, ymin=0, ymax=140, enlarge x limits=0.22,
  symbolic x coords={First rule, New rule (two frames), New rule (three frames)}, xtick=data,
  ylabel={Eye alarms per hour}, tick label style={font=\footnotesize}, label style={font=\footnotesize},
  ymajorgrids, grid style={gray!20}, axis x line*=bottom, axis y line*=left,
  nodes near coords, nodes near coords style={font=\footnotesize, /pgf/number format/fixed, /pgf/number format/precision=1},
  legend cell align=left, legend style={font=\footnotesize, draw=none, at={(0.99,0.98)}, anchor=north east}]
  \addplot[fill=own1, draw=none] coordinates {(First rule,97.1) (New rule (two frames),22.7) (New rule (three frames),7.4)};
  \addlegendentry{Drivers 1--9 (tuned)}
  \addplot[fill=oth1, draw=none] coordinates {(First rule,120.3) (New rule (two frames),20.7) (New rule (three frames),9.7)};
  \addlegendentry{Drivers 10--14 (not tuned)}
  \draw[dashed, line width=0.8pt, gray!70!black] (axis cs:{[normalized]-1},3) -- (axis cs:{[normalized]3},3);
  \node[font=\scriptsize, text=gray!70!black, anchor=south east, yshift=1pt] at (axis cs:{[normalized]0.72},3) {target: 3 per hour};
\end{axis}
\end{tikzpicture}

\caption{Eye alarms per hour on 14 alert DMD drivers, before and after the change}
\label{fig:eye-rate}
\end{figure}

How many real closures the new rule misses remains unknown, as we have no labelled closed-eye footage, and real closures recorded at full frame rate on the C922 are needed before these values can be trusted.

\subsection{Where the time goes}

We timed each stage on the replay clip to estimate what a C++ rewrite would gain (Table \ref{tab:stages}). Most of the Python only coordinates native libraries, so a full rewrite would save very little, and the detector cost does not apply on the Pi, where the \gls{npu} takes over.

The one every-frame cost under our control was the blind check, which therefore became a small C++ library. It takes 0.32 ms against 3.0 ms for OpenCV and reaches the same decision on all 15{,}723 DMD frames, with the OpenCV version kept as a fallback when the library is not built.

\begin{table}[htbp]
\centering
\caption{Time per frame by stage on the PC CPU}
\label{tab:stages}
\footnotesize
\setlength{\tabcolsep}{3pt}
\renewcommand{\arraystretch}{1.15}
\begin{tabular}{|p{0.330\linewidth}|p{0.230\linewidth}|p{0.360\linewidth}|}
\hline
\textbf{Stage} & \textbf{Time} & \textbf{Where the work runs} \\
\hline
Detector, CPU path only & 58.6 ms & Mostly native, 45.5 ms is the model \\
\hline
Eye module & About 11 ms & MediaPipe, native \\
\hline
Blind check & 3.0 $\rightarrow$ 0.32 ms & Moved to C++ \\
\hline
Video decode & 1.2 ms & OpenCV, native \\
\hline
\gls{ema} and decision & About 0.002 ms & Python \\
\hline
\end{tabular}
\end{table}

\section{Status of every specification}

Table \ref{tab:fs-status} gives the state of each specification today. ``Tested'' means the software behaviour was checked on a PC, and no entry implies that the specification has passed its final test on the Pi.

{\footnotesize
\setlength{\tabcolsep}{3pt}
\renewcommand{\arraystretch}{1.12}
\begin{longtable}{|p{0.075\linewidth}|p{0.250\linewidth}|p{0.130\linewidth}|p{0.470\linewidth}|}
\caption{Specification status at the end of SDP~1}
\label{tab:fs-status}\\
\hline
\textbf{Spec} & \textbf{Feature} & \textbf{State} & \textbf{Evidence and what is left} \\
\hline
\endfirsthead
\hline
\textbf{Spec} & \textbf{Feature} & \textbf{State} & \textbf{Evidence and what is left} \\
\hline
\endhead
\hline
\endfoot
FS-01 & Pi 5 deployment & Needs hardware & Setup, service and kiosk files written, never run on a Pi \\ \hline
FS-02 & Camera & Partly tested & Tested on video files, the live C922 path is not run \\ \hline
FS-03 & Detector & Partly tested & Runs with the DMD model, needs training on our subjects \\ \hline
FS-04 & \gls{ema} & Tested & T-04 passes on synthetic scores \\ \hline
FS-05 & Light and voice & Partly tested & Flash limit and clip selection pass, needs the real bar and speaker \\ \hline
FS-06 & Cool-down & Tested & T-06 passes, but see the phone-call case above \\ \hline
FS-07 & Mirror checks & Needs people & T-07 needs scripted mirror checks \\ \hline
FS-08 & Logging & Tested & Replay events stored and read back \\ \hline
FS-09 & On-device video & Tested & Upload contains event fields only \\ \hline
FS-10 & Dataset & Pilot done & Four unlabelled pilot sessions with three of the authors (Section \ref{sec:pilot}), the full dataset is still to be recorded \\ \hline
FS-11 & Pruning and quantization & Partly tested & Pruning, fine-tuning, \gls{onnx} and NCNN ran, Hailo, \gls{qat}, FP16 and \gls{int} did not \\ \hline
FS-12 & Hard negatives & Needs data & Miner ran but found none on a small sample, so it has yet to produce output \\ \hline
FS-13 & Baseline clips & Needs camera & Recording tool written, not run \\ \hline
FS-14 & Class alerts & Tested & Each class selects its own sound and pattern \\ \hline
FS-15 & Tiers & Tested & T-15 passes, and exact on the replay \\ \hline
FS-16 & Explained alerts & Tested & T-16 passes for every class \\ \hline
FS-17 & Blind check & Partly tested & T-17 passes in software, limits need tuning on the C922 \\ \hline
FS-18 & Crop tracking & Partly tested & Follows a real face, stays off until T-18 shows it beats the fixed crop \\ \hline
FS-19 & Calibration & Tested & Formula and redo rule pass, false-alarm benefit needs our subjects \\ \hline
FS-20 & Eye closure & Known issue & Timing passes, but false alarms are above target (Section \ref{sec:eye-results}) \\ \hline
FS-21 & Power-cut logging & Partly tested & 50 process kills clean, relay test still needed \\ \hline
FS-22 & Dashboard and sync & Tested & Real \gls{tls} upload to the dashboard on SQLite, PostgreSQL and hosting not run \\ \hline
FS-23 & Tips & Tested & Rules checked against known histories \\ \hline
FS-24 & Tablet & Partly tested & Screens checked in a browser, needs the 7-inch display and users \\ \hline
FS-25 & Pre-trip check & Tested & Logic and rest advice pass, user trial open \\ \hline
FS-26 & Rest stops & Partly tested & T-26 passed on the saved list, live query not proven, phone location needs HTTPS on the Pi \\ \hline
\end{longtable}
}


\section{What cannot be claimed yet}

The results above show that the software works as designed on a PC, but not yet that the product meets its targets. Real \gls{fps}, latency, heat and power on the Pi are unmeasured, as is everything that needs the AI HAT+, the light bar, the speaker or the C922. Per-class recall and false alarms per hour on our own subjects require the dataset and the leave-one-subject-out test, while T-07, T-18, T-24 and all user trials require the demonstration rig of Section \ref{sec:demo}. Table \ref{tab:ch5-08} therefore stays empty until SDP~2. The four pilot sessions of Section \ref{sec:pilot} show real cabins, mounts and light, but they are neither labelled nor run through the model, so no detection result comes from them yet.
